\documentclass[11pt]{article}
\usepackage[letterpaper,left=1.02in,right=1.02in,top=.91in,bottom=.9in]{geometry}
\usepackage{amsmath,amssymb,amsthm,booktabs,xcolor,fancyhdr,hyperref}
\definecolor{ink}{RGB}{22,51,72}
\definecolor{accent}{RGB}{12,96,91}
\hypersetup{colorlinks=true,linkcolor=accent,citecolor=accent,urlcolor=accent,
  pdftitle={When blurred Dirichlet zeros determine their sources},pdfauthor={OpenAI ChatGPT}}
\pagestyle{fancy}\fancyhf{}\fancyhead[L]{\small\itshape Blurred Dirichlet zeros}
\fancyhead[R]{\small Recovery and indistinguishability}\fancyfoot[C]{\thepage}
\setlength{\headheight}{14pt}
\newtheorem{theorem}{Theorem}
\newtheorem{lemma}[theorem]{Lemma}
\newtheorem{corollary}[theorem]{Corollary}
\newtheorem{proposition}[theorem]{Proposition}
\newcommand{\RR}{\mathbb R}
\newcommand{\CC}{\mathbb C}
\newcommand{\dd}{\,\mathrm d}
\newcommand{\A}{\mathcal A}
\newcommand{\D}{\mathcal D}
\newcommand{\cX}{\mathcal X}
\newcommand{\cM}{\mathcal M}
\newcommand{\abs}[1]{\left|#1\right|}
\newcommand{\Real}{\operatorname{Re}}
\newcommand{\supp}{\operatorname{supp}}
\newcommand{\sech}{\operatorname{sech}}
\newcommand{\srcbase}{https://github.com/VladimirReshetnikov/ProveIt/blob/2a9f8e3ee74a1824d57546f0b1b4f1cef5df93d6/Analysis/GeneralizedRiemannHypothesis}
\setlength{\parskip}{2.7pt plus .4pt}
\setlength{\parindent}{1.1em}
\setlength{\emergencystretch}{1.6em}
\allowdisplaybreaks
\begin{document}
\thispagestyle{plain}
\noindent{\small\sffamily\color{accent} ANALYTIC NUMBER THEORY \quad / \quad INVERSE PROBLEMS}\par
\vspace{.65em}
{\noindent\LARGE\bfseries\color{ink} When blurred Dirichlet zeros determine\\[2pt] their sources\par}
\vspace{.38em}
{\noindent\large A sharp exponential-moment threshold and a universal probability kernel\par}
\vspace{.45em}
\noindent{\small November 27, 2026}\par
\vspace{.5em}
\noindent\textbf{Abstract.}
Each member of a finite family of primitive Dirichlet \(L\)-functions has its zero ordinates translated according to an unknown finite measure. Given only the sum of the translated zero measures, we prove that every translating measure is determined when all the measures have finite moments at a common pair of one-sided exponential rates whose sum is at least \(\pi\). For symmetric tails the critical exponent is \(\pi/2\). Conversely, a single smooth, strictly positive, even probability density with every strictly subcritical exponential moment makes the zero measures of \emph{all} primitive characters of the same conductor and parity identical after convolution. We determine exactly which finite families remain identifiable below the threshold, including families containing \(\zeta\). We also give a finite-prime interpolation formula and a recovery theorem from separated windows at large height. The zero-location input is the vendored Dirichlet GRH theorem. The proofs require neither simple zeros nor restrictions on coincidences of ordinates.

\section{The observation and the threshold}
For a primitive character \(\chi\) of conductor \(q_\chi\), write
\(Z_\chi=\sum_\gamma m_\chi(\gamma)\delta_\gamma\), where
\(\tfrac12+i\gamma\) runs over its distinct nontrivial zeros and
\(m_\chi(\gamma)\) is their multiplicity. Both signs of the ordinates retain their orientation.
The primitive character of conductor one is denoted by \(\zeta\).
For a finite family \(\cX\) of distinct primitive characters we observe
\begin{equation}\label{eq:obs}
             M=\sum_{\chi\in\cX} Z_\chi*\nu_\chi .
\end{equation}
An atom of \(\nu_\chi\) at \(u\) translates that entire zero list by \(u\).
A probability law gives the averaged translated list. Thus \eqref{eq:obs}
also describes the intensity obtained by applying the same marginal displacement law to each zero of a given character. The data in the theorem are this exact measure, with its real coordinate fixed.

For \(a,b>0\), let \(\cM_{a,b}\) consist of finite complex Borel measures with
\[
 \|\nu\|_{a,b}:=\int_{u<0}e^{-au}\dd|\nu|(u)+
                    \int_{u\ge0}e^{bu}\dd|\nu|(u)<\infty .
\]
Write \(\cM_d=\cM_{d,d}\). The standard bound
\(Z_\chi([v,v+1])\ll_\chi \log(3+|v|)\) \cite{count} and integration in \(u\)
give the same bound for \(|M|\), with a constant depending on the input. In particular, \(M\) is locally finite and tempered.

\begin{theorem}[Sharp recovery threshold]\label{th:main}
For \(a+b\ge\pi\), the map \eqref{eq:obs} is injective on \(\cM_{a,b}^{\cX}\) for every finite family \(\cX\).
For \(0<a+b<\pi\), it is injective exactly in the following cases:
a singleton, a pair of nonprincipal characters of opposite parity, or a pair consisting of \(\zeta\) and one nonprincipal character. The empty family is trivial. In every other case distinct tuples of smooth strictly positive probability densities give the same \(M\).
\end{theorem}

\noindent At the critical width, even the set of characters can be unknown in advance: apply the theorem to the union of two candidate finite families, filling missing entries with zero. This identifies the characters that have nonzero measures. A lower bound on the separation of shifts would exclude continuous distributions and is unnecessary here.

\begin{theorem}[One probability law erases all prime frequencies]\label{th:universal}
There is a smooth, strictly positive, even probability density \(K\) on \(\RR\) with
\(\int e^{d|u|}K(u)\dd u<\infty\) for every \(d<\pi/2\), such that
\[
 Z_{\chi}*K=Z_{\chi'}*K
 \quad\Longleftrightarrow\quad
 q_\chi=q_{\chi'}\ \text{and}\ \chi(-1)=\chi'(-1).
\]
In particular the same \(K\) works for all conductors simultaneously. For a finite multiset of primitive characters blurred by this law, the complete combined observation determines precisely the total number of factors, the number of odd factors, the number of zeta factors and the product of the conductors.
\end{theorem}

For example, let \(\chi\) be the quartic character modulo \(5\) with
\(\chi(2)=i\). Its conjugate has \(\overline\chi(2)=-i\). Both have conductor \(5\) and take the value \(-1\) at \(4\). Theorem~\ref{th:universal} gives equality of their entire blurred zero measures. Theorem~\ref{th:main} rules out such an equality for \emph{any} probability law with a finite critical exponential moment. The distinction between one endpoint moment and all strictly smaller moments is therefore necessary.

\section{What the explicit formula measures}\label{sec:explicit}
We take \(\widehat f(x)=\int f(t)e^{-itx}\dd t\), extend this convention by duality to distributions, and put \(N_\chi(z)=\int e^{-izu}\dd\nu_\chi(u)\).
Then \(N_\chi\) is bounded and holomorphic on \(-a<\Im z<b\). Its derivatives of all orders are bounded on the real line, so multiplication by \(N_\chi\) preserves Schwartz tests. Fubini's theorem, the unit-interval bound, and the moments justify
\(\widehat M=\sum_\chi N_\chi\widehat Z_\chi\).
Write \(\epsilon_\chi=0\) or \(1\) according as \(\chi(-1)=1\) or \(-1\), and let \(\psi=\Gamma'/\Gamma\).

For an entire test \(h\) with smooth compactly supported Fourier transform the Guinand--Weil formula is
\begin{align}
2\pi\langle Z_\chi,h\rangle
={}&2\pi{\bf1}_{q_\chi=1}\{h(i/2)+h(-i/2)\}
+\int h(t)\{\log(q_\chi/\pi)+g_{\epsilon_\chi}(t)\}\dd t \notag\\
&-\sum_{n\ge2}\frac{\Lambda(n)}{\sqrt n}
  \{\chi(n)\widehat h(\log n)+\overline{\chi(n)}\widehat h(-\log n)\},
\label{eq:gw}\\[-4pt]
&\hspace{28mm}g_\epsilon(t)=\Real\psi\big((1/2+\epsilon+it)/2\big).\notag
\end{align}
Here \(\Lambda(p^k)=\log p\) for primes \(p\), and \(\Lambda(n)=0\) otherwise.
Equation~\eqref{eq:gw} uses the angular-frequency convention of the Article's equation \texttt{rec:eq:gwstrip}. It is the classical explicit formula \cite{gw,brs}. Dirichlet GRH is used to replace the complex arguments \((\rho-1/2)/i\) on its left by real ordinates.

\begin{lemma}[The atoms and the continuous density]\label{lem:spec}
At \(x=\log n\) and \(-\log n\), \(n=p^k\), the respective Dirac coefficients of \(\widehat Z_\chi\) are
\[
 -\frac{\Lambda(n)}{\sqrt n}\overline{\chi(n)}
 \quad\text{and}\quad -\frac{\Lambda(n)}{\sqrt n}\chi(n).
\]
Off zero and these frequencies, its density is the even function
\begin{equation}\label{eq:arch}
 D_\chi(x)=-\frac{e^{-(1/2+\epsilon_\chi)|x|}}{1-e^{-2|x|}}
              +{\bf1}_{q_\chi=1}\,2\cosh(x/2).
\end{equation}
The conductor term in the full transform is \(\log(q_\chi/\pi)\delta_0\).
\end{lemma}
\begin{proof}
Insert \(h=\widehat\varphi\) for \(\varphi\in C_c^\infty(\RR)\) in \eqref{eq:gw}. Then \(\widehat h(v)=2\pi\varphi(-v)\), which fixes the conjugations and all constants. The digamma series for \(g_\epsilon\) has, up to constants in \(t\), the terms
\(-2d_n/(d_n^2+t^2)\), \(d_n=2n+\epsilon+1/2\), \(n\ge0\).
Their Fourier transforms are \(-2\pi e^{-d_n|x|}\). The partial sums of the regularized digamma series are \(O(\log(3+|t|))\) uniformly in the cutoff: split at \(n=1+|t|\) and bound the tail by \(O(n^{-2}+t^2n^{-3})\). They therefore converge as tempered distributions. Off zero the exponential series converges locally uniformly and sums to \eqref{eq:arch} after division by \(2\pi\). The pole terms give the two exponentials making up \(2\cosh(x/2)\).
\end{proof}
The decomposition is local on bounded frequency intervals. The separate prime and pole terms can grow exponentially near infinity. The full identity is applied to compactly supported frequency tests.

\section{Why the critical width is \texorpdfstring{\(\pi\)}{pi}}\label{sec:width}
Let \(Q\) be a common multiple of the conductors, \(G=(\mathbb Z/Q\mathbb Z)^\times\), and inflate the characters to \(G\). They remain distinct, since a Dirichlet character has a unique primitive inducing character. For \(r\in G\) define the analytic function
\begin{equation}\label{eq:row}
 A_r(z)=\sum_{\chi\in\cX}\overline{\chi(r)}N_\chi(z),\qquad
 N_\chi(z)=\frac1{|G|}\sum_{r\in G}\chi(r) A_r(z).
\end{equation}
For \(Q=1\) the group is trivial. The second equality is orthogonality.
By Lemma~\ref{lem:spec}, the atom at \(+\log p\), divided by \(-\log p/\sqrt p\), equals \(A_r(\log p)\) when \(p\equiv r\pmod Q\). The atom at \(-\log p\) gives \(A_r(-\log p)\) when \(p\equiv r^{-1}\pmod Q\). Both readings apply to complex measures.

\begin{lemma}[Prime logarithms determine a bounded strip function]\label{lem:strip}
Put \(W=a+b\ge\pi\). If \(F\) is bounded and holomorphic on \(-a<\Im z<b\), and \(F(\log p)=0\) for all primes in one reduced residue class modulo \(Q\), then \(F=0\).
\end{lemma}
\begin{proof}
Let \(c=(b-a)/2\). The change of variable \(w=\exp(\pi(z-ic)/W)\) maps the strip onto the right half-plane, which the map \((w-1)/(w+1)\) sends onto the unit disk. A zero at \(\log p\) maps first to \(w_p=p^{\pi/W}e^{-i\pi c/W}\). Since \(|c|<W/2\),
\[
 1-\left|\frac{w_p-1}{w_p+1}\right|^2
       =\frac{4\Real w_p}{|w_p+1|^2}\asymp_{a,b} p^{-\pi/W}.
\]
The zeros of a nonzero bounded analytic disk function satisfy the Blaschke condition \(\sum(1-|z_j|)<\infty\), by Jensen's formula after removing any zero at the origin. The sum of \(1/p\) in each reduced residue class diverges \cite{mertens}. For \(\pi/W\le1\) this contradicts the Blaschke condition.
\end{proof}

Apply the lemma to the difference of the two candidate rows in \eqref{eq:row}. The equality of the observations gives the required zeros. Orthogonality and uniqueness of the Fourier transform of a finite measure prove the first assertion of Theorem~\ref{th:main}. All members of the family use the same strip. In the symmetric critical strip the disk coordinate is simply \(\tanh(z/2)\), and the deficit at a positive prime is exactly \(2/(p+1)\). The prime number distribution enters through this divergent series.

Two variants follow at once. One can omit any primes as long as the retained reciprocal-prime sum diverges in each required class at critical width. Also, any family of laws supported on a fixed half-line with one positive exponential moment in the unbounded direction has unique recovery: the exponent on the bounded side can be chosen arbitrarily large.

\section{What survives below the threshold}\label{sec:below}
Here \(a,b>0\) can be arbitrarily small. Take a difference of two tuples with the same \(M\), and use \(N_\chi\) for the corresponding difference transforms. Put
\[
 E=\sum_{\epsilon_\chi=0}N_\chi,\quad
 O=\sum_{\epsilon_\chi=1}N_\chi,\quad
 Z=N_\zeta\ \text{if present, and }0\text{ otherwise},\quad H(z)=2+2\cosh z.
\]
The sum \(E\) includes \(Z\). The continuous part of Lemma~\ref{lem:spec} implies for \(x>0\)
\[
 E(x)+e^{-x}O(x)=P(x)Z(x),\qquad
 P(z)=(e^z+1)(1-e^{-2z}).
\]
The corresponding equation on \(x<0\) uses \(e^x\) and \(P(-x)\).
Initially both identities hold between the prime frequencies. By analytic continuation each holds on the whole strip. Subtracting and using
\(P(z)=H(z)(1-e^{-z})\) gives the two identities
\begin{equation}\label{eq:par}
                    E=HZ,\qquad O=-HZ .
\end{equation}
Without \(\zeta\), the two parity sums of the translating measures are therefore uniquely determined. If at most one character of each parity occurs, all the individual measures are determined. With \(\zeta\) and just one other character, either \(O=0\) or \(E=Z\). In both cases \eqref{eq:par} also forces \(Z=0\). This proves the exceptional injectivity statements.

For the converse let \(W=a+b<\pi\), \(c=(b-a)/2\). The following entire function vanishes at \(0\) to order two and at \(\pm\log n\) for every integer \(n\ge2\):
\begin{equation}\label{eq:F}
 F_c(z)=e^{-(z-ic)\sinh(z-ic)}\sin(\pi e^z)\sin(\pi e^{-z}).
\end{equation}
It is nonzero, as evaluation at \(z=\log(3/2)\) shows. On every horizontal closed substrip \(|\Im z-c|\le d<\pi/2\) its modulus decays faster than \(e^{-L|\Real z|}\) for every \(L>0\). Indeed, writing \(v=\Im z-c\), \(x=\Real z\),
\[
 \Real((x+iv)\sinh(x+iv))=x\sinh x\cos v-v\cosh x\sin v .
\]
The first term eventually exceeds \(C_d|x|e^{|x|}\) with \(C_d>0\), whereas the other term and the logarithmic growth of the sines are \(O_d(e^{|x|})\). Also \(F_c(-\bar z)=\overline{F_c(z)}\).
Contour shifting in the inverse Fourier integral now gives a smooth real density with transform \(F_c\), with finite \((a,b)\) norm. For the right tail shift the contour strictly above \(b\), and for the left tail strictly below \(-a\). The same argument holds after multiplication by \(H\), any polynomial in \(\cosh z\), or \(z^{2k}\).

If two nonprincipal characters have the same parity, assign \(F_c,-F_c\) to their two difference transforms and zero to the others. If \(\zeta\), a nonprincipal even character and an odd character are present, assign them, in that order,
\begin{equation}\label{eq:triple}
                        F_c,\quad (H-1)F_c,\quad -H F_c .
\end{equation}
These possibilities include every family outside the exceptional cases. In either construction \eqref{eq:par} holds identically, every prime-power atom vanishes, and every conductor term vanishes at zero. To verify the full distribution at the origin, observe that \(g_\epsilon''\) is integrable by Stirling's formula. Hence \(x^2\widehat g_\epsilon=-\widehat{g_\epsilon''}\) is a continuous function. All the multipliers just constructed are smoothly divisible by \(x^2\). Their products with the gamma distributions are continuous across zero, so cancellation away from zero extends across it. The pole density is already smooth. Thus the \emph{complete} observation is zero for this nonzero difference tuple.

Let \(f_j\) be these real inverse transforms, with \(\int f_j=0\). For each \(j\) normalize \(g_j(t)=\sqrt{f_j(t)^2+e^{-2t^2}}\) to a probability density \(b_j=g_j/\int g_j\). For one sufficiently small \(\lambda>0\), all \(b_j+\lambda f_j\) and \(b_j-\lambda f_j\) are strictly positive smooth probability densities. They give two distinct admissible tuples with the same observation. This completes Theorem~\ref{th:main}. In the symmetric case \(c=0\), the densities are even and have every strictly subcritical moment at once. Multiplying \eqref{eq:F} by \(z^{2k}\) also makes any prescribed finite number of ordinary moments agree between the two tuples.

\section{A common positive law for all conductors}\label{sec:positive}
The signed construction already gives a loss of identifiability when the laws are unknown. Theorem~\ref{th:universal} is stronger: a single law can erase a distinction even for an observer who knows the law exactly. We construct its Fourier transform to vanish at all \(\pm\log n\), \(n\ge2\), while keeping its value at zero equal to one and its inverse transform positive.

\begin{lemma}[Making a signed density positive by convolution]\label{lem:pos}
Let \(A>0\). Suppose \(h\) is real, \(\int |h(s)|e^{d|s|}\dd s<\infty\) for every \(0\le d<A\), and its bilateral Laplace transform \(L(y)=\int h(s)e^{ys}\dd s\) satisfies \(\inf_{|y|<A}L(y)=m>0\). There is an even, integrable \(k>0\) with the same subcritical moments such that
\((h*k)(t)\ge (m/2)k(t)\) everywhere.
\end{lemma}
\begin{proof}
Choose \(0=d_0<d_1<\cdots\uparrow A\) and increasing \(L_j\ge j\) with
\(\int_{|s|>L_j}|h(s)|e^{d_j|s|}\dd s<m/4\).
Choose decreasing \(\eta_j>0\) so small that
\((\eta_j L_j^2/2)\int |h(s)|e^{d_j|s|}\dd s<m/4\).
Define an even convex \(C^1\) function \(V\) with \(V(0)=V'(0)=0\).
On the positive half-line, in stage \(j\) increase \(V'\) linearly from \(d_{j-1}\) to \(d_j\) at rate \(\eta_j\), then hold it equal to \(d_j\) for length \(L_{j+1}\). These stages cover the half-line since their plateau lengths are at least one.

Put \(k=e^{-V}\). When \(|t|\) lies in stage \(j\), \(|V'(t)|\le d_j\), and on \([t-L_j,t+L_j]\) the derivative is \(\eta_j\)-Lipschitz. The preceding constant interval has length \(L_j\), and subsequent ramps have slope at most \(\eta_j\). In stage one the global Lipschitz bound suffices. Convexity gives
\(k(t-s)/k(t)\le e^{sV'(t)}\). For \(|s|\le L_j\), Taylor's inequality and \(1-e^{-v}\le v\) for \(v\ge0\) bound the difference by \(e^{d_j|s|}\eta_j L_j^2/2\). For \(|s|>L_j\) the difference is bounded by \(e^{d_j|s|}\). Integration against \(|h(s)|\dd s\) yields
\[
  \left|\frac{(h*k)(t)}{k(t)}-L(V'(t))\right|<m/2.
\]
Finally \(V(t)/|t|\to A\) by averaging the slopes, so \(k\) and \(h*k\) have every strictly subcritical absolute exponential moment.
\end{proof}

Set \(S(w)=\sin(\pi w)/(1-w)\), using the removable value \(S(1)=\pi\), and define
\[
 U_c(z)=e^{-(z-ic)\sinh(z-ic)}S(e^z)S(e^{-z}).
\]
The estimates for \eqref{eq:F} hold for \(U_c\) as well. For \(c=0\) its inverse transform \(h\) is real, even and smooth. It and all its derivatives have every exponential moment of rate less than \(\pi/2\), by the same contour shifts. Its Laplace transform on the open interval is
\[
                U_0(iy)=e^{y\sin y}\,|S(e^{iy})|^2 .
\]
The right-hand side extends to a strictly positive continuous function on \([-\pi/2,\pi/2]\): the only integer on this half of the unit circle is \(1\), where \(S(1)=\pi\). Lemma~\ref{lem:pos} gives \(h*k>0\). Normalize it to integral one and call the resulting density \(K\). It is even and smooth by differentiation under convolution. Its Fourier transform is a constant multiple of \(U_0\widehat k\), so
\begin{equation}\label{eq:kill}
                \widehat K(0)=1,\qquad
                \widehat K(\pm\log n)=0\quad(n=2,3,\ldots).
\end{equation}

To read the information left by \eqref{eq:kill}, subtract two packets, and let \(e,o,z_0\) be the differences in their counts of even, odd and zeta factors (the even count includes zeta). If the observations agree then \eqref{eq:par} gives \(e=H(x)z_0\), \(o=-H(x)z_0\) near zero, since \(\widehat K(0)=1\).
The function \(H\) is nonconstant, so \(e=o=z_0=0\).
Both gamma distributions then cancel globally, as do the pole densities. What remains is a mass at zero equal to the logarithm of the ratio of the conductor products. Its vanishing proves necessity. Conversely, matching these four invariants cancels the gamma, pole and conductor contributions, and \eqref{eq:kill} cancels every arithmetic atom. Testing on arbitrary compactly supported smooth frequency functions proves equality of the full observations. The singleton case is exactly the conductor-and-parity equivalence in Theorem~\ref{th:universal}.

There is an asymmetric version. Fix \(|c|<\pi/2\), take \(h_c\) with transform \(U_c\), and apply Lemma~\ref{lem:pos} to \(e^{ct}h_c(t)\) with \(A=\pi/2\). Its Laplace transform at \(y\) is \(U_c(i(c+y))\), positive with a positive infimum by continuity on \(|y|\le\pi/2\): the angles \(c+y\) lie strictly between \(-\pi\) and \(\pi\). Untilting the positive convolution gives \(h_c*(e^{-ct}k)\). Its normalized law still has \eqref{eq:kill}, and has finite bilateral Laplace moments at every real parameter in \((c-\pi/2,c+\pi/2)\). Thus any prescribed subcritical asymmetric moment class also contains a common positive law with the same erasure property.

\section{Finite prime data and the endpoint}\label{sec:finite}
At the symmetric endpoint put \(A=\pi/2\). Consider a single row \(A_r\) in \eqref{eq:row}, bounded on the strip by \(M_0=\sum_\chi\|\nu_\chi\|_{A,A}\). For a difference of tuples one uses its corresponding bound. At \(t\in\RR\) put \(w=\tanh(t/2)\). For the primes up to \(X\), form the two sets of real disk points
\[
 s_j=\frac{p-1}{p+1}\quad(p\equiv r\bmod Q),\qquad
 s_j=-\frac{p-1}{p+1}\quad(p\equiv r^{-1}\bmod Q).
\]
These are all distinct. Write \(B(w)=\prod_j(w-s_j)/(1-s_jw)\) and \(D_j=A_r(2\operatorname{arctanh}s_j)\), which is measured by the corresponding prime atom.

\begin{proposition}[An explicit convergent reconstruction]\label{pr:decode}
The rational interpolant
\[
       I_X(w)=\sum_j D_j\,\frac{B(w)}{B'(s_j)(w-s_j)}
\]
takes its values at the nodes by continuity. For real \(t\),
\begin{align}
 \abs{A_r(t)-I_X(w)}
   &\le M_0\cosh(t/2)\,\abs{B(w)},                                      \label{eq:err}\\
 |B(w)|&=
 \prod_{\substack{p\le X\\p\equiv r\ (Q)}}\left|\frac{p-e^t}{p+e^t}\right|
 \prod_{\substack{p\le X\\p\equiv r^{-1}\ (Q)}}\left|\frac{p-e^{-t}}{p+e^{-t}}\right|
 \ll_{Q,t} (\log X)^{-4\cosh(t)/\varphi(Q)} .                           \label{eq:rate}
\end{align}
At fixed \(Q,t\), the implicit constant may be taken uniformly in \(r\). With errors of size \(\delta_j\) in the input atoms \emph{after} division by \(-\log p/\sqrt p\), add
\(\sum_j\delta_j|B(w)/(B'(s_j)(w-s_j))|\) to the right of \eqref{eq:err}.
\end{proposition}
\begin{proof}
Set \(f(w)=A_r(2\operatorname{arctanh}w)\). In the unit disk it is bounded by \(M_0\). Integrate \(f(v)/[B(v)(v-w)]\) on a circle of radius tending to one. The residue formula gives \((f(w)-I_X(w))/B(w)\). The minimum of \(|B|\) on these circles tends to one. By Cauchy--Schwarz the limiting mean of \(|e^{i\theta}-w|^{-1}\) is at most \(1/\sqrt{1-|w|^2}\), proving \eqref{eq:err}. If the row vanishes at the nodes one also has the simpler bound \(|A_r(t)|\le M_0 |B(w)|\) by dividing out the disk factors and applying the maximum principle.

The elementary formula for each disk factor gives the products. Away from a vanishing factor, expand \(\log|(p-u)/(p+u)|=-2u/p+O_u(p^{-3})\) for \(p>2u\), first with \(u=e^t\) and then with \(u=e^{-t}\). Mertens' theorem in progressions \cite{mertens} gives \eqref{eq:rate}. The error term for inexact readings follows by linearity of \(I_X\). For fixed \(t\), applying \eqref{eq:row} to all reconstructed rows recovers each \(N_\chi(t)\).
\end{proof}

For the quartic example, use \(Q=5\), \(r=2\), \(r^{-1}=3\). At \(t=0\), the product in \eqref{eq:rate} is \(0.06579295\) for \(X=100\) and \(0.02623577\) for \(X=100000\). At \(t=\log(3/2)\) it is respectively \(0.03074662\) and \(0.01131906\). These are finite-product illustrations. Very close disk nodes can amplify errors in individual atom measurements, as the explicit noise term records.

There is also a finite cost for \emph{exact} erasure. If a probability law \(K_*\)
has critical moment \(M_*=\int e^{A|u|}\dd K_*(u)\), and
\(\widehat K_*(\pm\log n)=0\) for \(n=2,\ldots,L\), divide its transform by the corresponding disk factors and evaluate at \(w=0\). Since \(\widehat K_*(0)=1\), this gives
\[
   M_*\ge\prod_{n=2}^{L}\left(\frac{n+1}{n-1}\right)^2
       =\left(\frac{L(L+1)}2\right)^2 .
\]
The same product over a smaller set gives the necessary cost for erasing only those frequencies.

\begin{corollary}[Polynomial losses at the endpoint]\label{cor:end}
Injectivity still holds if, for each measure and some finite \(k\ge0\),
\[
       \int e^{A|u|}(1+|u|)^{-k}\dd|\nu_\chi|(u)<\infty,\qquad A=\pi/2.
\]
\end{corollary}
\begin{proof}
For \(\beta<A\) close to \(A\), a difference row has supremum \(O((A-\beta)^{-k})\) on \(|\Im z|<\beta\), with \(k\) chosen as the largest exponent among the finitely many inputs. Put \(s=\pi/(2\beta)>1\). A zero at \(\log p\) contributes the real disk factor \(|(p^s-e^{st})/(p^s+e^{st})|\) at \(t\). For fixed \(t>0\), take all \(p>2e^t\) in class \(r\). Since \(\log((1-v)/(1+v))\le-2v\) for \(0<v<1\), finite products followed by a limit give
\[
 |A_r(t)|\ll (A-\beta)^{-k}
       \exp\!\left[-2e^{st}\sum_{\substack{p>2e^t\\p\equiv r\ (Q)}}p^{-s}\right].
\]
Partial summation of the reciprocal-prime asymptotic gives
\(\sum_{p\equiv r}p^{-s}=\varphi(Q)^{-1}\log(1/(s-1))+O_Q(1)\) as \(s\downarrow1\). Indeed this is its Laplace integral against
\((s-1)e^{-(s-1)v}\dd v\) after writing the partial sum through \(e^v\) as
\(\varphi(Q)^{-1}\log v+O(1)\) for \(v\ge2\).
Since \(s-1\asymp A-\beta\), the bound tends to zero whenever \(2e^t/\varphi(Q)>k\). Thus the row vanishes on a ray and hence everywhere. Apply \eqref{eq:row}.
\end{proof}
In particular the single probability density of Theorem~\ref{th:universal} fails \emph{every} polynomially discounted critical moment, although all smaller exponential moments are finite. Otherwise the two different quartic characters modulo \(5\), each bearing that law, would contradict Corollary~\ref{cor:end}.

\section{Recovery from sparse high windows}\label{sec:windows}
\begin{proposition}[Windowed observations]\label{pr:win}
Let $M=\sum_\chi\nu_\chi*Z_\chi$ denote the observed measure and
$\widehat\nu(v)=\int e^{-ivy}\,d\nu(y)$. Fix any $T_j\to+\infty$.
The restrictions of $M$ to $[T_j-H_j,T_j+H_j]$ determine the two normalized prime readings at $\pm\log p$ for every prime $p$ if
\[
 H_j-\pi^{-1}\log\log\log T_j\longrightarrow+\infty.
\]
Here it suffices that $\int\log(3+|y|)\,d|\nu_\chi|(y)<\infty$.
Consequently the injective conclusion for \(a+b\ge\pi\) in Theorem~\ref{th:main} holds from these restrictions alone. The observation also allows an error measure on each window whose
total variation there tends to zero.

\end{proposition}
\begin{proof} Fix a prime $p$, put $u=\log p$, and fix $m\ge4$.
Use the Article's lemmas on decay on the closed strip and frequency response
in its subsection on characters and multiplicity vectors in short windows.
For $Q\ge2$ bounding all conductors and $T=T_j$, put
\[
 \ell=\log(e^eQ(T+3)),\quad R=32\log\ell,\quad
 D=2+\pi^{-1}\log(4R),\quad b=8p,\quad B=bm.
\]
Set $K=J_b^{*m}*k_R$, where $J_b=(2b)^{-1}\mathbf{1}_{[-b,b]}$ and
$k_R(x)$ is the Article's normalized positive kernel proportional to
$\exp[-\cosh(\pi x/2)-R\cosh(\pi x)/(R+\cosh(\pi x))]$.
The test $h(t)=(iu)^{-1}(e^{iut}K(t-T))'$ obeys
$\widehat h(0)=0$ and $\widehat h(u)=1$.
In the explicit formula with test $h(t+y)$, translation multiplies each coefficient
$\widehat h(v)$ by $e^{ivy}$. The kernel is holomorphic across the closed critical strip and decays rapidly there, so these shifted tests are admissible in the explicit formula. The frequency-response lemma bounds the sum
of the absolute values of all unwanted prime-power terms, after multiplication
by $\sqrt p/u$, by $C2^{-m}$ uniformly in $T$ and $y$.

For the remaining terms, the conductor constant integrates to zero.
Integration by parts bounds the gamma integral in absolute value by
$u^{-1}\int K(s-T)|G'_\chi(s-y)|\,ds$, where
$G_\chi(s)=\log(q_\chi/\pi)+\Re\psi((1/2+a_\chi+is)/2)$,
$a_\chi$ is the parity and $\psi=\Gamma'/\Gamma$.
For each fixed $y$ the integral tends to zero because $G'_\chi$ is bounded and tends to
zero at infinity, and the probability densities $K$ are uniformly tight
about zero for fixed $p,m$ by the strip lemma.
The bound is uniform in $y,T$. For the zeta factor the pole evaluations
$h(y\pm i/2)$ are likewise uniformly bounded and tend to zero for each fixed $y$,
by the same lemma applied on the two edges of the strip.
Bounded convergence therefore handles integration against each finite $\nu_\chi$.

The unit-interval zero count gives
$|M|([v,v+1])\ll_{M,Q}\log(3+|v|)$ by the logarithmic moment assumption.
For $|t-T|\ge D+2B$ the strip lemma gives
$|h(t)|\ll_{p,m} e^{-R/4}\exp[-(|t-T|-B)]$.
Summing over unit intervals, the total discarded integral is
$O_{M,Q,p,m}(e^{-R/4}(1+\log(3+T)))=o(1)$.
Since $D=\pi^{-1}\log\log\log T+O_Q(1)$, the observed radius eventually
exceeds $D+2B$ for each fixed $p,m$. Thus the observable
\[
 E_{p,m,j}=-\frac{2\pi\sqrt p}{u}
       \int_{T_j-H_j}^{T_j+H_j}h(t)\,dM(t)
\]
satisfies
$\limsup_{j\to\infty}|E_{p,m,j}-\sum_\chi\chi(p)\widehat\nu_\chi(-u)|
\le C2^{-m}\sum_\chi\|\nu_\chi\|_{\mathrm{TV}}$.
Taking $m\to\infty$ determines the coefficient uniquely.
Using $\overline h$ as the test gives
$\sum_\chi\overline{\chi(p)}\widehat\nu_\chi(u)$ by the identical estimate.
For fixed $p,m$ the tests are uniformly bounded on the real line, which
also proves the stated error tolerance.\end{proof}

The centers \(T_j\) may grow arbitrarily rapidly. The additive excess in the half-width lets one choose a fixed prime and a fixed accuracy parameter before sending \(j\) to infinity. This distinguishes the infinite-dimensional problem from the Article's single-window identification of a bounded finite packet. No finite observation precision or bounded list of primes provides the exact infinite-data conclusion by itself.

\section{Source relationship and further questions}
The pinned ProveIt revision is \texttt{2a9f8e3ee74a1824d57546f0b1b4f1cef5df93d6}. Its vendored \texttt{grh/} tree matches upstream revision \texttt{72fead102f663d5c7a407964a13dc0cf9d977207} \cite{lean}. In \texttt{Statements.lean}, the primitive zero-location assertion quantifies over all primitive characters and all zeros in \(0<\Real s<1\). In \texttt{Main.lean}, the theorem \texttt{LFunction\_eq\_zero\_re\_of\_isPrimitive} supplies precisely the specialization used here, and \texttt{riemannHypothesis} supplies the zeta case. The upstream proof derives half-plane nonvanishing from a one-sided estimate for primes in progressions. Our argument uses its zero-location output. The Fourier and probability arguments in this paper are separate mathematical proofs. The Lean declarations were inspected at the pinned revision. The build and independent proof-checker runs described in its README are upstream checks. This paper has no Lean formalization.

The Article's README describes a merger of 120 reports with 869 result statements \cite{article}. Its recovery section establishes single-window identification of finite products, rank bounds for finitely many translated lists, and character recovery from balanced statistics of pooled ordinates. The nearest antecedents here are its explicit formula, the translated-family theorem \texttt{rec:thm:moving}, and the localized test-function lemmas \texttt{rec:lem:lekernel} and \texttt{rec:lem:lefreq}. The translations in that theorem range over compact configurations of finitely many distinct character-shift pairs. Here a single character can carry a continuous law, its support can be unbounded, and only the summed observation is available.

The explicit formula, the strip zero-set criterion and character orthogonality are classical. Bondarenko, Radchenko and Seip construct interpolation bases using zeros of \(L\)-functions \cite{brs}. Ruzsa and Sz\'ekely study positive convolution factors for signed functions of positive integral \cite{rs}. The moment control up to an open strip boundary in Lemma~\ref{lem:pos} allows the resulting density to attain the exact arithmetic threshold. A targeted comparison of the Article and the literature located these antecedents. The proposed new results are the threshold for unknown displacement laws, the positive density erasing all prime-power readings with all subcritical moments, and the corresponding classification. Historical priority beyond that comparison is open.

\paragraph{Further questions.}
What is the optimal modulus of continuity for recovering the translating measures when only finitely many noisy atoms or high windows are observed? Can the unbounded additive excess in Proposition~\ref{pr:win} be replaced by a fixed constant for continuous displacement laws? Which more precise critical weights, between polynomial discounts and arbitrary subcritical exponential moments, still force identification? For other families of \(L\)-functions, the method also asks which independent gamma and pole profiles survive after a positive law removes the arithmetic atoms.

\begingroup\small\setlength{\parskip}{0pt}\begin{thebibliography}{9}\setlength{\itemsep}{2pt}\setlength{\parsep}{0pt}
\bibitem{lean} L. Sernau repository, \emph{The Generalized Riemann Hypothesis for Dirichlet \(L\)-functions, in Lean 4}. Vendored readable edition, credited there to the proof of Levent Alp\"oge with Claude. \href{\srcbase/vendor/grh-lean/README.md}{Import record},
\href{\srcbase/vendor/grh-lean/grh/src/Statements.lean}{statement definitions},
\href{\srcbase/vendor/grh-lean/grh/src/Main.lean}{exported theorems}.
\bibitem{article} V. Reshetnikov, ProveIt repository, \emph{Beyond the Riemann hypothesis for Dirichlet \(L\)-functions: consequences, transfers and limits}, merged article. \href{\srcbase/Article/README.md}{Overview} and \href{\srcbase/Article/sections/rec.tex}{section on recovering zeros and characters}.
\bibitem{gw} A. Chirre, V. J. Pereira J\'unior and D. de Laat, \emph{Primes in arithmetic progressions and semidefinite programming}, Math. Comp. \textbf{90} (2021), 2235--2246, equations (2.1) and (2.6). \href{https://doi.org/10.1090/mcom/3638}{doi:10.1090/mcom/3638}.
\bibitem{count} M. A. Bennett, G. Martin, K. O'Bryant and A. Rechnitzer, \emph{Counting zeros of Dirichlet \(L\)-functions}, Math. Comp. (2021). \href{https://doi.org/10.1090/mcom/3599}{doi:10.1090/mcom/3599}.
\bibitem{mertens} K. S. Williams, \emph{Mertens' theorem for arithmetic progressions}, J. Number Theory \textbf{6} (1974), 353--359.
\href{https://doi.org/10.1016/0022-314X(74)90032-8}{doi:10.1016/0022-314X(74)90032-8}. For the reciprocal-prime version see A. Languasco and A. Zaccagnini, \emph{Computing the Mertens and Meissel--Mertens constants for sums over arithmetic progressions}, Exp. Math. \textbf{19} (2010), 279--284, \href{https://arxiv.org/abs/0906.2132}{arXiv:0906.2132}.
\bibitem{brs} A. Bondarenko, D. Radchenko and K. Seip, \emph{Fourier interpolation with zeros of zeta and \(L\)-functions}, Constr. Approx. \href{https://doi.org/10.1007/s00365-022-09599-w}{doi:10.1007/s00365-022-09599-w}.
\bibitem{rs} I. Z. Ruzsa and G. J. Sz\'ekely, \emph{Convolution quotients of nonnegative functions}, Monatsh. Math. \textbf{95} (1983), 235--239. \href{https://doi.org/10.1007/BF01352002}{doi:10.1007/BF01352002}.
\end{thebibliography}\endgroup
\end{document}
